% !TEX root =  main.tex
\chapter{Introduction}
\section{Background of Study}
\hspace{\parindent}


\section{Statement of the Problem}
\hspace{\parindent} 


\section{Objectives of the Study}
\hspace{\parindent} The study aims to investigate the effects of coherence in the performance of the phase retrieval algorithm specifically, the Single-Beam Multiple-Intensity Reconstruction (SBMIR). Through understanding the principles of speckle contrast and formation, and coherence, a model for the use of partially coherent illumination in SBMIR will be developed. Phase, as well as the amplitude, reconstructions will be analyzed throughout the investigation. The model will utilize a recently proposed enhanced SBMIR method that involves non-sequential propagation which addresses some issues of the conventional one. 

\section{Significance of the Study}
\hspace{\parindent}For years, partially coherent beams (PCB) have gained considerable attention because of their useful applications on different fields of optics and photonics. Compared to highly coherent beams, illumination from a partially coherent source reduces speckle effects and improves the image quality \cite{rayleigh}. These partially coherent light sources are found to be applied on free-space optical communications \cite{2020,chang}. Implementation using a partially coherent light source could extend the utilization and application of phase retrieval.

By introducing PCB light sources to the phase retrieval algorithm, specifically the intensity recording part, the configuration cost can be reduced significantly since they are more insensitive to external perturbations. The ability to maximize the available intensity measurements and information will make the proposed technique more cost-effective. Moreover, numerical simulations can provide valuable insights on a viable model for phase retrieval using partially coherent light.

\section{Scope and Limitation of the Study}
\hspace{\parindent}The reconstruction algorithms to be analyzed in this study are applied on numerically simulated intensity images. Reconstructions from experimental data are not in the scope of the study. In addition, the main basis for the OMW method is the Fresnel diffraction integral which was consequently used to simulate wave propagation. However, the Fresnel diffraction integral is known to scale with propagation distance and wavelength which limits the parameters that can be used in the numerical simulations [6]. Lastly, phase retrieval methods involving multiple intensity measurements are highly dependent on the intensity variations. Consequently, complete wave reconstruction cannot be always guaranteed as the reconstruction algorithm is object dependent as well.